
\subsection*{2.3.2}
\phantomsection
\label{sga3:VIIB:2.3.2}
\addcontentsline{toc}{subsection}{\SGARefLink{sga3:VIIB:2.3.2:prop}{2.3.2}}

现在为简便起见，假设环 \(k\) 为 Artin 环。当 \(G\) 在 \(k\) 上拓扑平坦时，可以用 Cartier 的一个泛性质刻画代数 \[
  \mathbf H(G)=\mathbf H(G)(k)
\]。回顾 \textup{（参见 \SGARefLink{sga3:VIIA:1.2.1}{1.2.1}）}，若 \(U\) 为 \(k\)-模，则以 \(\mathbf W(U)\) 表示将每个有限长度 \(k\)-代数 \(C\) 映为 \(C\)-模 \[
  U\otimes_k C.
\] 的函子。\footnote{\label{sga3:VIIB:nde:88}编者补充了前一句，并在下文用 \(\mathbf W(U)^\times\) 代替 \(U_m\)。还须回顾 \textup{（参见 \SGARefLink{sga3:VIIA:1.2.1}{1.2.1}）}，“\(k\)-函子”是从 \(\operatorname{Alf}/k\) 到 \((\mathbf{Sets})\) 的协变函子。}若 \(U\) 为 \(k\)-代数 \textup{（结合且有单位元）}，则 \(U\otimes_k C\) 亦然。以 \(\mathbf W(U)^\times\) 表示将每个 \(C\in\operatorname{Ob}(\operatorname{Alf}/k)\) 映为代数 \(U\otimes_k C\) 的可逆元乘法群的 \(k\)-群函子：\[
  \mathbf W(U)^\times(C)=(U\otimes_k C)^\times.
\]。

此外，将 \(G\) 与 \(k\)-群函子 \[
  C\longmapsto
  \operatorname{Hom}_{\operatorname{Vaf}/k}
    \bigl(\operatorname{Spf}(C),G\bigr),
\] 等同，并以 \[
  \operatorname{Hom}_{k\text{-}\operatorname{Gr.}}
    \bigl(G,\mathbf W(U)^\times\bigr)
\] 表示从 \(G\) 到 \(\mathbf W(U)^\times\) 的 \(k\)-群函子同态的集合。有下述命题。

\medskip\noindent\textbf{命题。}\phantomsection \label{sga3:VIIB:2.3.2:prop}{\itshape 设 \(k\) 为 Artin 环。对每个在 \(k\) 上拓扑平坦的形式群 \(G\) 及每个 \(k\)-代数 \(U\)，存在典范同构 \[
  \operatorname{Hom}_{k\text{-}\operatorname{Gr.}}
    \bigl(G,\mathbf W(U)^\times\bigr)
  \xrightarrow{\ \sim\ }
  \operatorname{Hom}_{k\text{-}\operatorname{Alg.}}
    \bigl(\mathbf H(G),U\bigr).
\]。}

设 \(A\) 为 \(G\) 的仿射代数。由假设，\(A\) 为 \(\operatorname{PC}(k)\) 的投射对象，且 \[
  \mathbf H(G)=A^\dagger.
\]。对 \(\operatorname{PC}(k)\) 的每个对象 \(P\)，以 \[
  U\widehat\otimes_k P
\] 表示各 \(k\)-模 \[
  U\otimes_k(P/N),
\] 的逆极限，其中 \(N\) 遍历 \(P\) 的开子模。存在线性映射 \[
  U\otimes_k(P/N)
  \longrightarrow
  \operatorname{Hom}_k\bigl((P/N)^*,U\bigr)
\]，将 \(u\otimes x\) 映为 \(k\)-线性映射 \[
  f\longmapsto f(x)u
\]，并构成滤过逆系。因此取逆极限得到态射 \begin{equation}
  U\widehat\otimes_kP
  \xrightarrow{\ \psi_P\ }
  \operatorname{Hom}_k(P^\dagger,U).
  \tag{1}
  \label{sga3:VIIB:2.3.2:eq1}
\end{equation}。

当 \(P=k\) 时，态射 \(\psi_k\) 显然是同构。此外，\textup{\SGARefLink{sga3:VIIB:2.3.2:eq1}{(1)}} 两端视为 \(P\) 的函子时，均与无限乘积交换 \textup{（每个乘积都是有限乘积的滤过逆极限）}。因此当 \(P\) 为若干个 \(k\) 的副本的乘积时，\textup{\SGARefLink{sga3:VIIB:2.3.2:eq1}{(1)}} 为同构；进而当 \(P\) 为 \(\operatorname{PC}(k)\) 的投射对象时亦然 \textup{（\textup{\SGARefLink{sga3:VIIB:2.3.2:eq1}{(1)}} 两端均与有限直和交换）}。

现在以 \[
  \operatorname{Hom}_{\mathcal F}\bigl(G,\mathbf W(U)\bigr)
\] 表示从 \(G\) 到 \(\mathbf W(U)\) 的 \(k\)-函子态射的集合。由于 \[
  G=\operatorname{Spf}(A)
   =\varinjlim_{\mathfrak l}\operatorname{Spf}(A/\mathfrak l),
\]，其中 \(\mathfrak l\) 遍历 \(A\) 的开理想，有典范同构 \[
\begin{aligned}
  \operatorname{Hom}_{\mathcal F}\bigl(G,\mathbf W(U)\bigr)
  &=
  \varprojlim_{\mathfrak l}
  \operatorname{Hom}_{\mathcal F}
    \bigl(\operatorname{Spf}(A/\mathfrak l),\mathbf W(U)\bigr)
  \\
  &=
  \varprojlim_{\mathfrak l}
  U\otimes_k(A/\mathfrak l)
  =
  U\widehat\otimes_kA.
\end{aligned}
\]。因此前面的论证给出典范同构 \begin{equation}
  \operatorname{Hom}_{\mathcal F}\bigl(G,\mathbf W(U)\bigr)
  =
  U\widehat\otimes_kA
  \xrightarrow{\ \psi_A\ }
  \operatorname{Hom}_k\bigl(\mathbf H(G),U\bigr).
  \tag{2}
  \label{sga3:VIIB:2.3.2:eq2}
\end{equation}。

对每个有限长度 \(k\)-代数 \(C\)，乘法使 \(U\otimes_kC\) 成为含幺幺半群，而每个含幺幺半群同态 \[
  G(C)\longrightarrow U\otimes_kC
\] 必然是群同态 \[
  G(C)\longrightarrow(U\otimes_kC)^\times.
\]。因此 \[
  \operatorname{Hom}_{k\text{-}\operatorname{Gr.}}
    \bigl(G,\mathbf W(U)^\times\bigr)
\] 是 \(\operatorname{Hom}_{\mathcal F}(G,\mathbf W(U))\) 中由取值于含幺幺半群的 \(k\)-函子态射构成的部分。

还须说明，这些态射对应于从 \(\mathbf H(G)\) 到 \(U\) 的保持乘法与单位元的 \(k\)-线性映射。\footnote{\label{sga3:VIIB:nde:89}原文称，这“由 \(\psi_A\) 的函子性推出”。编者补充了下面的细节。}为简化记号，写作 \[
  H=\mathbf H(G)=A^\dagger
\]，并以 \(\widehat\otimes\) 代替 \(\widehat\otimes_k\)。以 \[
  \Delta_A,\ m_A,\ \varepsilon_A
  \qquad\text{且}\qquad
  \Delta_H,\ m_H,\ \varepsilon_H
\] 分别表示 \(A\) 和 \(H\) 的余乘法、乘法与增广。设 \[
  \theta\in
  \operatorname{Hom}_{\mathcal F}\bigl(G,\mathbf W(U)\bigr),
\]，\(\gamma\) 为其在 \(U\widehat\otimes A\) 中的像，\[
  \phi:H\longrightarrow U
\] 为对应的 \(k\)-线性映射。态射 \(\theta\) 将单位截面 \(s\in G(k)\) 映为 \(U\) 的元素 \(u\)。由于 \(s\) 对应于增广 \[
  \varepsilon:A\longrightarrow k,
\]，即 \(H\) 的单位元 \(1_H\)，有 \[
  \theta(s)=1_U
  \quad\Longleftrightarrow\quad
  \phi(1_H)=1_U.
\]。

态射 \[
  \theta\circ m_G:G\times G\longrightarrow\mathbf W(U)
\] 对应于元素 \[
  (\operatorname{id}_U\widehat\otimes\Delta_A)(\gamma),
\]，而后者通过对偶对应于映射 \[
  \phi\circ m_H:H\otimes H\longrightarrow U.
\]。

另一方面，态射 \[
  \theta\circ\operatorname{pr}_1:
  G\times G\longrightarrow\mathbf W(U)
\] 对应于元素 \[
  \gamma\widehat\otimes1_A
  \in
  U\widehat\otimes A\widehat\otimes A,
\]，后者通过对偶对应于 \[
  \phi\circ(\operatorname{id}_H\otimes\varepsilon_H):
  H\otimes H\longrightarrow U.
\]。类似地，\(\theta\circ\operatorname{pr}_2\) 对应于元素 \[
  \tau(\gamma\widehat\otimes1_A)
  \in
  U\widehat\otimes A\widehat\otimes A,
\]，其中 \[
  \tau(u\widehat\otimes a\widehat\otimes b)
  =
  u\widehat\otimes b\widehat\otimes a,
\]；此元素通过对偶对应于 \[
  \phi\circ(\varepsilon_H\otimes\operatorname{id}_H):
  H\otimes H\longrightarrow U.
\]。

最后，由 \[
\begin{aligned}
  &(U\widehat\otimes A\widehat\otimes A)
   \times
   (U\widehat\otimes A\widehat\otimes A)
  \longrightarrow
  U\widehat\otimes A\widehat\otimes A,
  \\
  &\bigl(
    u\widehat\otimes a_1\widehat\otimes a_2,\,
    u'\widehat\otimes a_3\widehat\otimes a_4
  \bigr)
  \longmapsto
  uu'\widehat\otimes a_1a_3\widehat\otimes a_2a_4
\end{aligned}
\] 给出的乘法映射 \[
  \mu=m_{U\widehat\otimes A\widehat\otimes A}
\] 可视为 \[
  (U\otimes U)\widehat\otimes A^{\widehat\otimes4},
\] 的“交换因子 \(a_2\) 和 \(a_3\)”的自同态 \(\sigma_{23}\) 与映射 \[
  (U\otimes U)
  \widehat\otimes A^{\widehat\otimes2}
  \widehat\otimes A^{\widehat\otimes2}
  \xrightarrow{\ m_U\widehat\otimes m_A\widehat\otimes m_A\ }
  U\widehat\otimes A\widehat\otimes A.
\] 的复合。

由此可知，映射 \[
  \mu\circ
  \bigl(\theta\circ\operatorname{pr}_1,\,
        \theta\circ\operatorname{pr}_2\bigr):
  G\times G\longrightarrow\mathbf W(U)
\] 对应于下面的复合映射 \(\beta\)：\begin{center}\phantomsection \label{sga3:VIIB:2.3.2:beta-diagram} % SGA3-VII-DIAG-105
% Authority: Exp7B-13oct24.pdf, local p. 51.
% Crop: figures/Exp7B_localp051_cartier_universal_beta_diagram.png
% Locator: tex/components/41_expose_VIIB_section232_cartier_universal_property_en.tex, Loop 1 line 308.
\begin{tikzcd}[row sep=2.7em,column sep=3em]
{} & H\otimes H
     \arrow[r,"\beta"]
     \arrow[dl,"{\operatorname{id}_H\otimes\sigma_{23}\otimes
                   \operatorname{id}_H}"']
   & U \\
H\otimes H\otimes H\otimes H
  \arrow[dr,"{(\operatorname{id}_H\otimes\varepsilon_H)
              \otimes(\varepsilon_H\otimes\operatorname{id}_H)}"']
  & {} & U\otimes U \arrow[u,"m_U"'] \\
{} & H\otimes H \arrow[ur,"\phi\otimes\phi"'] & {.}
\end{tikzcd}
\end{center}

最后，\(\theta\) 与 \(G\) 和 \(\mathbf W(U)^\times\) 的运算相容，当且仅当 \[
  \theta\circ m_G
  =
  \mu\circ
  \bigl(\theta\circ\operatorname{pr}_1,\,
        \theta\circ\operatorname{pr}_2\bigr),
\]。\noindent\textit{原文说明。}\enspace 法文再版在这两个有序对中印作 \(\phi\circ\operatorname{pr}_i\)。由于 \(\operatorname{pr}_i:G\times G\to G\)，而 \(\phi:H\to U\)，规范记录 \texttt{SGA-SRC-000847} 恢复了类型正确的 \(\theta\circ\operatorname{pr}_i\)。

由前面的论证，这等价于 \[
  \phi\circ m_H=\beta.
\]。但显然 \[
  (\phi\circ m_H)(x\otimes y)=\phi(xy),
\]，而 \[
  \beta(x\otimes y)=\phi(x)\phi(y).
\]。
